\chapter{Introduction}

This document is intended to be both an example of the TU Delft \LaTeX{} template for reports which can be used for both BSc and MSc theses, as well as a short introduction to its use. It is not intended to be a general introduction to \LaTeX{} itself,\footnote{We recommend \url{http://en.wikibooks.org/wiki/LaTeX} as a reference and a starting point for new users.} and we will assume the reader to be familiar with the basics of creating and compiling documents.

Instructions on how to use this template under Windows, Linux and MacOS, and which \LaTeX{} packages are required, can be found in \texttt{README.txt}.


\section{Document Structure}

Since a report, and especially a thesis, might be a substantial document, it is convenient to break it up into smaller pieces. In this template we therefore give every chapter its own file. The chapters (and appendices) are gathered together in \texttt{report.tex}, which is the master file describing the overall structure of the document. \texttt{report.tex} starts with the line
\begin{verbatim}
  \documentclass{tudelft-report}
\end{verbatim}
which loads the TU Delft report template. The template is based on the \LaTeX{} \texttt{scrbook} document class and stored in \texttt{tudelft-report.cls}. The document class accepts several comma-separated options. The default language is (UK) English, but this can be changed to Dutch (\emph{e.g.}, for bachelor theses) by specifying the \texttt{dutch} option:
\begin{verbatim}
  \documentclass[dutch]{tudelft-report}
\end{verbatim}
Furthermore, hyperlinks are shown in blue, which is convenient when reader the report on a computer, but can be expensive when printing and does not serve the purpose it has in a digital document. They can be turned black with the \texttt{print} option. This will also turn the chapter numbers grey instead of light blue.

If the document becomes large, it is easy to miss warnings about the layout in the \LaTeX{} output. In order to locate problem areas, add the \texttt{draft} option to the \texttt{\textbackslash documentclass} line. This will display a vertical bar in the margins next to the paragraphs that require attention.

This template has the option to generate a cover page by including the file \texttt{cover.tex}. See the next section for a detailed description.

The contents of the report are included between the \texttt{\textbackslash{}begin\{document\}} and \texttt{\textbackslash{}end\{docu\-ment\}} commands, and split into four parts by
\begin{enumerate}
  \item \texttt{\textbackslash{}frontmatter}, which uses Roman numerals for the page numbers and is used for the title page and the table of contents;
  \item \texttt{\textbackslash{}mainmatter}, which uses Arabic numerals for the page numbers and is the style for the chapters;
  \item \texttt{\textbackslash{}appendix}, which uses letters for the chapter numbers, starting with `A'.
  \item \texttt{\textbackslash{}backmatter}, which has no chapter numbers at all.
\end{enumerate}
The title page is defined in a separate file, \texttt{title.tex}. Additionally, it is possible to include a preface, containing, for example, the acknowledgements. An example can be found in \texttt{preface.tex}. The table of contents is generated automatically with the \texttt{\textbackslash{}tableofcontents} command. Chapters are included after \texttt{\textbackslash{}mainmatter} and appendices after \texttt{\textbackslash{}appendix}. For example, \texttt{\textbackslash include\{in\-troduction\}} includes \texttt{introduction.tex}.

\subsection{Bibliography}

Lists of references are handled by \texttt{biblatex}.
The present document has the following preamble
\begin{verbatim}
\usepackage[style=apa]{biblatex}
\addbibresource{report.bib}
\end{verbatim}
The first line tells us that the references will be formatted according 
to the APA-standards.

The second line points to a file that contains the references in a database-like format.
Use one such line for every database that you want to use.

The general way this works is as follows (using the present document as an example):
\begin{verbatim}
lualatex report
biber report
lualatex report
lualatex report
\end{verbatim}
The first line typesets the document and gathers information for the lists of references; the second line will generate lists of references for all chapters, these will be read in during execution of the third and fourth lines.
The repetition of \texttt{lualatex report} is needed to get all the page numbers in links and references right.
Depending on the editor that you use much of this can be done by clicking on a suitable icon or by hitting suitably defined hot keys. On Overleaf, all of this is done when you press \texttt{ctrl+r} or \texttt{ctrl+enter}.

The \texttt{biblatex} package has many options and can be tailored to almost every need; you can explore its documentation at \texttt{https://www.ctan.org/pkg/biblatex}

To get numbered references just use
\begin{verbatim}
\usepackage{biblatex}
\addbibresource{report.bib}
\end{verbatim}
That is, remove \texttt{style=apa}.

To give some examples we cite an article: \textcite{Einstein1906},
a book: \textcite{MR1039321}, another article: \textcite{MR3860876},
and a book with more than one editor: \textcite{MR3204729}.


\section{Cover and Title Page}

This template will automatically generate a cover page if you include the file \texttt{cover.tex}.
This file defines three commands:

\texttt{\textbackslash{}MakeCoverFront} inserts just the front cover

\texttt{\textbackslash{}MakeCoverFull} inserts the complete cover (front, back, and spine) with some `bleed', i.e., a slightly enlarged edge to allow for cutting of the page after printing.

\texttt{\textbackslash{}MakeCoverBack} inserts just the back cover

\texttt{cover.tex} can be adapted according to your liking; this process is relatively free-format, with some examples given in \texttt{cover.tex}.

For a thesis it is desirable to have a title page within the document, containing information like the thesis committee members. To give you greater flexibility over the layout of this page, it is described in the file \texttt{title.tex}. Modify this file according to your needs. The example text is in English, but Dutch translations are provided in the comments. Note that for a thesis, the title page is subject to requirements which differ by faculty. Make sure to check these requirements before printing.


\section{Chapters}

Each chapter has its own file. For example, the \LaTeX{} source of this chapter can be found in \texttt{introduc\-tion.tex}. A chapter starts with the command
\begin{verbatim}
  \chapter{Chapter title}
\end{verbatim}
This starts a new page, prints the chapter number and title and adds a link in the table of contents. If the title is very long, it may be desirable to use a shorter version in the page headers and the table of contents. This can be achieved by specifying the short title in brackets:
\begin{verbatim}
  \chapter[Short title]{Very long title with many words 
    which could not possibly fit on one line}
\end{verbatim}

Chapters are subdivided into sections, subsections, subsubsections, and, optionally, paragraphs and subparagraphs. All can have a title, but only sections and subsections are numbered:
\section{\textbackslash section\{\ldots\}}
\subsection{\textbackslash subsection\{\ldots\}}
\subsubsection{\textbackslash subsubsection\{\ldots\}}
\paragraph{\textbackslash paragraph\{\ldots\}}

\section{Fonts and Colors}

The main font in this template is Arial, with the main headers in Roboto Slab, according to the TU Delft corporate style.
Alternatively, using the class option \texttt{\textbackslash{}documentclass[fourier]}, the serif font Utopia can be used (provided by the package \texttt{fourier}).
In that latter case, Helvetica will be used in the main headers.

The corporate colors of the TU Delft are \textcolor{tud Blue}{\textbf{blue}}, \textbf{black} and \textbf{white}, available, respectively, via \texttt{\textbackslash text\-color\{tud Blue\}\{...\}}, \texttt{\textbackslash textcolor\{black\}\linebreak[0]\{...\}} and \texttt{\textbackslash textcolor\{white\}\linebreak[0]\{...\}}.
Apart from these three, the house style defines the following basic colors, that can be used via \texttt{\textbackslash text\-color\{tud \textit{colorname}\}\{...\}}, where \texttt{\textit{colorname}} is any of:

\vspace{\topsep}\noindent
\begin{minipage}{.32\textwidth}
  \begin{itemize}
    \itemsep 0pt
    \parskip 0pt
    \item \textcolor{tud DarkBlue}{\textbf{Dark blue}}
    \item \textcolor{tud Turquoise}{\textbf{Turquoise}}
    \item \textcolor{tud RoyalBlue}{\textbf{Royal blue}}
    \item \textcolor{tud LightPurple}{\textbf{Light purple}}
  \end{itemize}
\end{minipage}
\begin{minipage}{.32\textwidth}
  \begin{itemize}
    \itemsep 0pt
    \parskip 0pt
    \item \textcolor{tud Pink}{\textbf{Pink}}
    \item \textcolor{tud Burgundy}{\textbf{Burgundy}}
    \item \textcolor{tud Red}{\textbf{Red}}
    \item \textcolor{tud Orange}{\textbf{Orange}}
  \end{itemize}
\end{minipage}
\begin{minipage}{.32\textwidth}
  \begin{itemize}
    \itemsep 0pt
    \parskip 0pt
    \item \textcolor{tud Yellow}{\textbf{Yellow}}
    \item \textcolor{tud Green}{\textbf{Green}}
    \item \textcolor{tud ForestGreen}{\textbf{Forest green}}
    \item \textcolor{tud DarkGrey}{\textbf{Dark grey}}
  \end{itemize}
\end{minipage}
\vspace{\topsep}

For shading of areas, specific reduced coverages of blue and dark grey may be used, but mind contrast and readability.
These reduced coverages are:
\begin{itemize} 
    \itemsep 0pt
    \parskip 0pt
  \item \textcolor{tud Blue}{\textbf{Blue:}}
    \tikz[baseline]{\node[fill=tud Blue!50, anchor=base] {50\,\%,};}%
    \tikz[baseline]{\node[fill=tud Blue!10, anchor=base] {10\,\%};}
  \item \textcolor{tud DarkGrey}{\textbf{Dark grey:}}
    \tikz[baseline]{\node[fill=tud DarkGrey!60, anchor=base] {\textcolor{white}{60\,\%,}};}%
    \tikz[baseline]{\node[fill=tud DarkGrey!40, anchor=base] {40\,\%,};}%
    \tikz[baseline]{\node[fill=tud DarkGrey!15, anchor=base] {15\,\%};}
\end{itemize}

Furthermore, there are paler versions of these colors to be used for data visualisation only:

\vspace{\topsep}\noindent
\begin{minipage}{.32\textwidth}
  \begin{itemize}
    \itemsep 0pt
    \parskip 0pt
    \item \textcolor{tud DarkBluePale}{\textbf{Dark blue}}
    \item \textcolor{tud TurquoisePale}{\textbf{Turquoise}}
    \item \textcolor{tud RoyalBluePale}{\textbf{Royal blue}}
    \item \textcolor{tud LightPurplePale}{\textbf{Light purple}}
  \end{itemize}
\end{minipage}
\begin{minipage}{.32\textwidth}
  \begin{itemize}
    \itemsep 0pt
    \parskip 0pt
    \item \textcolor{tud PinkPale}{\textbf{Pink}}
    \item \textcolor{tud BurgundyPale}{\textbf{Burgundy}}
    \item \textcolor{tud RedPale}{\textbf{Red}}
    \item \textcolor{tud OrangePale}{\textbf{Orange}}
  \end{itemize}
\end{minipage}
\begin{minipage}{.32\textwidth}
  \begin{itemize}
    \itemsep 0pt
    \parskip 0pt
    \item \textcolor{tud YellowPale}{\textbf{Yellow}}
    \item \textcolor{tud GreenPale}{\textbf{Green}}
    \item \textcolor{tud ForestGreenPale}{\textbf{Forest green}}
    \item \textcolor{tud DarkGreyPale}{\textbf{Dark grey}}
  \end{itemize}
\end{minipage}
